Com o objectivo de iniciar um estudo rigoroso de versões derivadas de teorias de campo quântico topológicas com defeitos, construímos uma $(\oo,1)$-categoria de bordismos estratificados e esboçamos também uma construção do seu análogo $(\oo,d)$. Na primeira parte da tese introduzimos um método de construção de $(\oo,1)$-categorias simétricas monoidais, inspirado nos espaços de Segal, mas mais adequado a construções geométricas, a que chamamos \emph{espaços de Segal intervalares}.

Usamos estes espaços para construir a $(\oo,1)$-categoria de variedades suaves de dimensão $d$, recorrendo à teoria de $\oo$-feixes suaves e corrigindo um erro no trabalho de Calaque--Scheimbauer. Dotamos depois esta estrutura de uma estrutura simétrica monoidal, calculamos a sua categoria de homotopia 1-truncada e construímos o espaço de Segal intervalar dos bordismos de dimensão $d$ com estruturas tangenciais.

Para captar as estratificações sobre as quais os defeitos topológicos são suportados, recorremos a \emph{espaços estratificados conicamente suaves}, uma versão adequadamente suave das estratificações cónicas. Esta escolha permite trabalhar simultaneamente com famílias de estratificações e com noções flexíveis de estrutura tangencial. Introduzimos então uma noção de \emph{variedade topológica estratificada} como $\cB$-variedade estratificada, onde $\cB$ é uma adequada $(\oo,1)$-categoria de básicos. Com esta classe de espaços construímos um espaço de Segal intervalar $\tilbordstr_d$ cujos objectos são variedades topológicas estratificadas compactas e fechadas de dimensão $(d-1)$ e cujos 1-morfismos são bordismos estratificados de dimensão $d$, dotando-o ainda de uma estrutura simétrica monoidal. Usamos também estruturas tangenciais para introduzir uma noção de rótulos de defeito.

Por fim, generalizamos $\Int$ para $\cube_d$ com $d>1$, delineando um modelo para $\Cat_{(\oo,d)}$, um $\cube_d$-espaço de bordismos estratificados e uma aplicação da homologia de factorização estratificada à construção de TQFTs estratificadas com defeitos.
\bigskip

\begin{samepage}
\noindent\textbf{Palavras-chave:}
\begin{itemize}[nosep]
    \item espaços estratificados conicamente suaves
    \item bordismos estratificados
    \item espaços de Segal intervalares
    \item categorias superiores
    \item teorias de campo quântico topológicas
\end{itemize}
\end{samepage}
